% Source: Wald GR, physical PDF p. 156
%         (printed p. 143).
% Coverage: Figure 6.5, effective potential for null geodesics.
% Figures/tables/footnotes: Figure 6.5; no tables or footnotes.
% Status: source-reviewed and compile-clean.
% Uncertainties: none.
\begingroup
\renewcommand{\thefigure}{6.5}
\begin{figure}[htbp]
\centering
\begin{tikzpicture}[x=.75cm,y=1.45cm,>=latex,font=\small]
  \draw[->] (0,0) -- (8.1,0) node[right] {\(r\)};
  \draw[->] (0,-1.45) -- (0,1.95) node[above] {\(V\)};
  \foreach \x/\lab in {1.25/\(3M\),2.1/\(5M\),4.25/\(10M\),6.4/\(15M\)}
    \draw (\x,.06) -- (\x,-.06) node[below=3pt] {\lab};
  \node[left] at (0,0) {0};
  \draw[thick,smooth] plot coordinates {
    (.82,-1.35) (.94,-.30) (1.13,.72) (1.36,1.12) (1.62,.93)
    (2.2,.62) (3.1,.34) (4.3,.18) (5.6,.10) (7.55,.045)};
\end{tikzpicture}
\caption{The effective potential, \(V\), for null geodesics.}
\label{fig:6.5}
\end{figure}
\endgroup
